% ============================================================================
% KEY MODULES
% ============================================================================

\section{Key Modules and Responsibilities}
\label{sec:modules}

This section details each major module in the \ctx{} codebase, its responsibilities, and key implementation details.

\subsection{CLI Layer}

\subsubsection{\file{cli.rs} (278 lines)}

Defines the command-line interface structure using the \code{clap} crate:

\begin{itemize}
    \item \code{Args} struct --- Global options (format, ignore patterns, verbosity)
    \item \code{Command} enum --- Top-level subcommands
    \item \code{QueryCommand} enum --- Query-specific subcommands
    \item \code{OutputFormat} enum --- Output format options (XML, Markdown, Plain, JSON)
\end{itemize}

\subsubsection{\file{main.rs} (1,272 lines, 43 symbols)}

The application entry point and command dispatcher:

\begin{itemize}
    \item Entry point with error handling
    \item Individual \func{run\_*} functions for each command
    \item Output formatting helpers
    \item Watch mode implementations
\end{itemize}

\subsection{Context Generation Pipeline}

\subsubsection{\file{walker.rs} (230 lines, 14 symbols)}

File discovery and filtering using the \code{ignore} crate:

\begin{lstlisting}[style=rustcode, caption={File walker configuration}]
let mut builder = WalkBuilder::new(root);
builder.git_ignore(true);           // Respect .gitignore
builder.add_custom_ignore_filename(".contextignore");

for entry in builder.build() {
    // Process file
}
\end{lstlisting}

Key features:
\begin{itemize}
    \item Gitignore support via the \code{ignore} crate
    \item Custom \code{.contextignore} file support
    \item Glob pattern matching
    \item Binary file detection
    \item 170+ built-in ignore patterns
\end{itemize}

\subsubsection{\file{tree.rs} (15 symbols)}

Project tree visualization with ASCII rendering:

\begin{verbatim}
project/
├── src/
│   ├── main.rs (1.2 KB)
│   └── lib.rs (856 B)
└── Cargo.toml (423 B)
\end{verbatim}

Features optional file size display and respects the ignore system.

\subsubsection{\file{formatter.rs} (47 symbols)}

Output formatting through a trait-based design:

\begin{lstlisting}[style=rustcode, caption={Formatter trait definition}]
pub trait Formatter {
    fn start(&self, tree_block: Option<&str>) -> String;
    fn format_file(&self, path: &str, content: &str) -> String;
    fn end(&self) -> String;
    
    // Streaming variants
    fn stream_start(&self, tree_block: Option<&str>) -> String;
    fn stream_file(&self, path: &str, content: &str) -> String;
    fn stream_end(&self) -> String;
}
\end{lstlisting}

Implementations:
\begin{itemize}
    \item \code{XmlFormatter} --- XML output with file paths as attributes
    \item \code{MarkdownFormatter} --- GitHub-flavored markdown with code blocks
    \item \code{PlainFormatter} --- Simple text format with separators
\end{itemize}

\subsubsection{\file{output.rs} (10 symbols)}

Orchestrates context generation:

\begin{itemize}
    \item \func{generate\_context()} --- Buffered output generation
    \item \func{stream\_context()} --- Streaming output for large codebases
\end{itemize}

\subsection{Code Intelligence Core}

\subsubsection{\file{index/mod.rs} (34 symbols)}

The indexing engine responsible for building the code intelligence database:

\begin{lstlisting}[style=rustcode, caption={Indexer structure}]
pub struct Indexer {
    db: Database,
    parser: CodeParser,
    root: PathBuf,
    verbose: bool,
}

impl Indexer {
    pub fn index_file(&mut self, path: &Path) -> Result<()>;
    pub fn index_all(&mut self) -> Result<IndexStats>;
    pub fn watch(&mut self) -> Result<()>;
}
\end{lstlisting}

Key responsibilities:
\begin{itemize}
    \item Incremental updates via content hashing
    \item Source compression with gzip (60-70\% reduction)
    \item Watch mode for auto-reindexing on file changes
    \item Batch processing for efficiency
\end{itemize}

\subsubsection{\file{db/schema.rs} (74 symbols)}

SQLite database operations:

\begin{itemize}
    \item Schema creation and migrations
    \item Symbol and edge CRUD operations
    \item FTS5 full-text search integration
    \item Duplicate detection with Jaccard similarity
    \item Embedding vector storage and retrieval
\end{itemize}

\subsubsection{\file{db/models.rs} (28 symbols)}

Data model definitions:

\begin{lstlisting}[style=rustcode, caption={Core data models}]
pub struct Symbol {
    pub id: String,
    pub name: String,
    pub kind: SymbolKind,
    pub file_path: String,
    pub line_start: u32,
    pub line_end: u32,
    pub visibility: Visibility,
    pub signature: Option<String>,
    pub brief: Option<String>,
    pub parent_id: Option<String>,
}

pub struct Edge {
    pub source_id: String,
    pub target_name: String,
    pub target_id: Option<String>,
    pub kind: EdgeKind,
    pub line: u32,
    pub context: Option<String>,
}

pub enum EdgeKind {
    Calls,
    Extends,
    Implements,
    Imports,
}
\end{lstlisting}

\subsubsection{\file{analytics/mod.rs} (31 symbols)}

DuckDB-powered analytical queries:

\begin{itemize}
    \item \func{get\_call\_graph()} --- Recursive call graph traversal
    \item \func{get\_callers()} --- Find all callers of a function
    \item \func{get\_impact()} --- Impact analysis for changes
    \item \func{analyze\_complexity()} --- Fan-in/fan-out complexity scoring
\end{itemize}

\subsection{Language Parsers}

All parsers reside in \file{src/parser/} and share common utilities from \file{mod.rs}.

\begin{table}[h]
\centering
\begin{tabular}{lrll}
\toprule
\textbf{Parser} & \textbf{Symbols} & \textbf{Languages} & \textbf{Key Extractions} \\
\midrule
\file{rust.rs} & 26 & Rust & fn, struct, enum, trait, impl \\
\file{typescript.rs} & 41 & TS/JS/TSX/JSX & function, class, interface, type \\
\file{python.rs} & 44 & Python & def, class, decorators, async \\
\file{solidity.rs} & 22 & Solidity & contract, function, event, error \\
\bottomrule
\end{tabular}
\caption{Language parser comparison}
\label{tab:parsers}
\end{table}

\subsubsection{Shared Utilities (\file{parser/mod.rs}, 40 symbols)}

Common functionality extracted to the module root:

\begin{itemize}
    \item \func{extract\_call\_edges()} --- Generic call extraction using tree-sitter queries
    \item \func{extract\_module\_name()} --- Module name derivation from file path
    \item \func{parse\_block\_doc\_comment()} --- Documentation comment parsing
    \item \code{SymbolKindMapping} --- Capture name to symbol kind mapping
\end{itemize}

\subsection{Embeddings System}

\subsubsection{\file{embeddings/mod.rs} (29 symbols)}

Core embedding infrastructure:

\begin{lstlisting}[style=rustcode, caption={Embedding provider trait}]
pub trait EmbeddingProvider {
    fn embed(&self, texts: &[String]) -> Result<Vec<Vec<f32>>>;
    fn dimension(&self) -> usize;
}

pub fn cosine_similarity(a: &[f32], b: &[f32]) -> f32 {
    let dot: f32 = a.iter().zip(b).map(|(x, y)| x * y).sum();
    let norm_a: f32 = a.iter().map(|x| x * x).sum::<f32>().sqrt();
    let norm_b: f32 = b.iter().map(|x| x * x).sum::<f32>().sqrt();
    dot / (norm_a * norm_b)
}
\end{lstlisting}

\subsubsection{\file{embeddings/local.rs} (21 symbols)}

Local embedding generation using \code{fastembed}:

\begin{itemize}
    \item Model: \code{all-MiniLM-L6-v2}
    \item Dimensions: 384
    \item Requirements: ~90MB download on first run
    \item No API key required
\end{itemize}

\subsubsection{\file{embeddings/openai.rs} (25 symbols)}

OpenAI API integration:

\begin{itemize}
    \item Model: \code{text-embedding-3-small}
    \item Dimensions: 1,536
    \item Requirements: \code{OPENAI\_API\_KEY} environment variable
    \item Batch processing with rate limiting
\end{itemize}
